\section{Tách kim loại khỏi hợp chất}

\subsection{Nguyên tắc chung}

Trong quặng, kim loại thường tồn tại dưới dạng hợp chất. Muốn thu kim loại phải chuyển ion kim loại về đơn chất. Phương pháp tách phụ thuộc mức độ hoạt động của kim loại và dạng hợp chất có trong quặng.

\begin{ghinho}{Ba quá trình cần nắm vững}
\textbf{Sắt từ iron(III) oxide:}
\[
\ce{Fe2O3 + 3CO ->[t^\circ] 2Fe + 3CO2}.
\]
\textbf{Nhôm từ aluminium oxide nóng chảy:}
\[
\ce{2Al2O3 -> 4Al + 3O2}\quad\text{(điện phân nóng chảy)}.
\]
\textbf{Kẽm từ zinc sulfide:}
\[
\ce{2ZnS + 3O2 ->[t^\circ] 2ZnO + 2SO2},
\]
\[
\ce{ZnO + C ->[t^\circ] Zn + CO}.
\]
\end{ghinho}

\subsection{Chọn phương pháp tách}

Kim loại rất hoạt động như Al không được tách khỏi oxide bằng CO ở mức kiến thức đang học; aluminium được sản xuất bằng điện phân nóng chảy \ce{Al2O3} trong hỗn hợp phù hợp. Với Fe, Zn, có thể dùng chất khử như CO hoặc C để khử oxide ở nhiệt độ cao.

\begin{vidu}
Vì sao không dùng CO để sản xuất Al từ \ce{Al2O3} theo cách sản xuất Fe từ \ce{Fe2O3}?
\end{vidu}
\begin{loigiai}
Al hoạt động hóa học mạnh hơn Fe và \ce{Al2O3} rất bền. Trong phạm vi chương trình, CO không được dùng để khử \ce{Al2O3} thành Al. Aluminium được sản xuất bằng điện phân nóng chảy aluminium oxide.

\textbf{Vậy,} phương pháp tách kim loại phải phù hợp với mức độ hoạt động của kim loại.
\end{loigiai}

\subsection{Đọc sơ đồ công nghệ}

Khi gặp một quy trình nhiều công đoạn, hãy ghi cho mỗi mũi tên ba thông tin:
\begin{enumerate}
\item chất đầu và chất cuối;
\item phương trình hoặc tỉ lệ mol;
\item hiệu suất nếu đề cho.
\end{enumerate}
Sau đó mới nối các công đoạn lại bằng số mol.

\begin{tuluyen}
Viết các phương trình hóa học cho quá trình tách Zn từ quặng chứa \ce{ZnS} theo hai giai đoạn: nung quặng trong không khí tạo \ce{ZnO}; sau đó khử \ce{ZnO} bằng carbon.
\dapso{$\ce{2ZnS + 3O2 -> 2ZnO + 2SO2}$; $\ce{ZnO + C -> Zn + CO}$.}
\end{tuluyen}
